\chapter{Introduction}\label{chap:introduction}

Modern financial markets run substantially on decision support systems and algorithmic trading. Institutional and retail participants alike depend on them to reduce large volumes of market information to something that can be acted upon \cite{hendershott_does_2011}. What that machinery consists of has changed considerably over two decades of research. Rule-based strategies came first \cite{brock1992}, classifiers drawn from machine learning followed \cite{gerlein2016,fischerkrauss2018}, and the most recent generation learns its trading policy from the data directly, by deep reinforcement learning \cite{huang2022,shavandi2022,hambly2023}.

Running alongside that progression, multi-agent systems (MAS) have been put forward as a way of imposing order on so many competing rules and models: independent agents each issue a signal, and a coordinator resolves those signals into one decision \cite{luo2002,korczak2013,hernes2016,hernes2024}. What recommends the arrangement is the separation it enforces between two questions that are easily run together, namely what a given indicator reports about the market and what a portfolio ought to do in response.

The practical problem that motivates this thesis is narrower than the breadth of that literature suggests. A retail investor on the Warsaw Stock Exchange (WSE) has access to daily price data, a modest brokerage account, and a large menu of published trading rules, but has no reliable way to tell which of those rules, or which combination of them, is worth following. The published evidence is not much help, because a substantial part of it is produced under evaluation protocols that make results look better than they are.

\section{Problem statement}\label{sec:problem}

Empirical studies of trading strategies are prone to four methodological weaknesses, and in combination these undermine the credibility of what such studies report.

\emph{In-sample parameter selection} is the first. Choosing the best configuration of a strategy on the very data that will afterwards be used to judge it inflates the apparent out-of-sample performance \cite{bailey2014,harvey2016}. \emph{Omission of realistic frictions} is the second, commissions and slippage above all, which between them can consume a noticeable share of whatever gross return a high-turnover configuration produces.

\emph{Neglect of risk-adjusted performance} is the third. Cumulative profit alone is a common basis for setting a trading strategy against passive buy-and-hold, and it leaves out of account that an active strategy may deliver a smoother equity curve while returning less unleveraged. Sustained bull markets sharpen the problem, since anything less than continuous exposure is mechanically outpaced there by a fully invested benchmark \cite{dichtl2020}. \emph{Absence of formal significance tests} is the fourth. Headline figures appear as bare point estimates, unaccompanied by confidence intervals or by paired tests against the benchmark, notwithstanding that the metrics involved are known to be non-Gaussian and path-dependent \cite{politis1994,bailey2014}.

A geographical concentration compounds these methodological ones. Evidence on multi-agent decision support systems clusters heavily on the major markets of the United States, Western Europe and Asia, whereas the Warsaw Stock Exchange, the largest in Central and Eastern Europe (CEE) by capitalisation, has drawn comparatively little study \cite{czapkiewicz2018,kompa2009,gurgul2025}. The difficulty facing a retail investor there is consequently twofold: the evidence available is methodologically fragile, and such of it as exists was collected somewhere else.

\section{Aim and scope of the thesis}\label{sec:aim}

The aim of this thesis is to design a transparent multi-agent decision support system for trading large-capitalisation stocks on the Warsaw Stock Exchange, and to determine, under an evaluation protocol that addresses the four weaknesses stated above, whether such a system offers a measurable advantage over passive investment.

The scope is defined by four decisions. The system trades daily closing prices of WIG20 index constituents, so intraday behaviour, order-book dynamics, and other asset classes are outside its remit. The signal agents use only price-derived technical indicators, so fundamental and alternative data are excluded. The aggregation rule is a linear, confidence-weighted score, chosen so that the decision layer remains interpretable and its mapping to forecast-combination theory stays transparent, in contrast to opaque end-to-end learners. Finally, all evidence is produced by simulation on historical data, so no live-trading exercise is part of the work.

Three hypotheses, stated formally in Section~\ref{sec:hypotheses} and tested in Chapter~\ref{chap:evaluation}, give the aim its testable form. In summary, they assert that the system matches or exceeds passive buy-and-hold on risk-adjusted metrics even where it does not on gross return (H1), that it outperforms the best single-agent baseline on total return (H2), and that its performance is robust to the removal of any single agent (H3).

The technical assumptions adopted throughout are the following. Positions are opened and closed at the daily closing price. Transaction costs comprise a proportional commission and one-way slippage, calibrated per market. Only decision-agent parameters are tuned, and only on a training window that precedes the evaluation window chronologically; all indicator hyperparameters are fixed at their literature-conventional values. Portfolio capital is allocated as a fixed fraction of available cash per position, without leverage.

\section{Own contribution}\label{sec:own_contribution}

The work reported in this thesis is the author's own. The system, the experiments, and the analysis were designed, implemented, and carried out by the author. Specifically, the author's contribution comprises the following:

\begin{itemize}
    \item the design and implementation of the complete software system described in Chapter~\ref{chap:implementation}, comprising five signal agents, the decision agent, the portfolio simulator, the metrics calculator, and the experiment harness,
    \item the design of the evaluation protocol of Chapter~\ref{chap:evaluation}, including the anchored walk-forward procedure, the transaction-cost calibration, the market-regime labelling, the volatility-matched comparison, and the bootstrap testing scheme,
    \item the implementation and in-sample tuning of the two machine-learning baselines used as reference points,
    \item the conceptual grounding of the aggregation rule in forecast-combination theory, together with the correlation and confidence-validation experiments that test how far that grounding holds empirically,
    \item the execution and interpretation of all experiments reported in Chapter~\ref{chap:evaluation}, including the external validation on the FTSE~100.
\end{itemize}

\section{Structure of the thesis}\label{sec:structure}

The thesis is organised into six chapters, whose significance is unequal by design: Chapters~\ref{chap:implementation} and~\ref{chap:evaluation} carry the original work, while the remaining chapters frame it.

Chapter~\ref{sec:related_work} reviews the state of the art across four streams: technical analysis and algorithmic trading; multi-agent systems in finance; the algorithmic trading literature on Central and Eastern European markets; and the evaluation of risk-adjusted performance. It establishes both the methodological gap and the empirical niche that the thesis addresses.

Chapter~\ref{sec:theory} presents the theoretical background and the system design. It introduces forecast-combination theory and the variance-reduction argument that motivates the architecture, maps that theory onto the trading problem, states the three hypotheses, and then describes the four-layer architecture, the five signal agents, the decision agent, and the agent taxonomy.

Chapter~\ref{chap:implementation} describes the software realisation: the requirements it had to satisfy, the technologies selected and why, the structure of the code, the data preparation pipeline, and the way the implementation was verified.

Chapter~\ref{chap:evaluation} reports the empirical part of the work. It describes the evaluation protocol in full, reports the results on the WIG20 panel, subjects them to a battery of robustness checks and an external validation on the FTSE~100, and interprets the evidence against the three hypotheses.

Chapter~\ref{sec:conclusion} summarises what was achieved, states the conclusions as an organised list, discusses the limitations of the study and the difficulties encountered during the work, and indicates directions for further development.
